%--- iliad ---
% title: Lecture 2: the correspondence argument
%--- end ---
% Lecture 2. Standalone Beamer source. No external assets required.
\documentclass[aspectratio=169,12pt]{beamer}
\usepackage[T1]{fontenc}
\usepackage{lmodern}
\usepackage{amsmath,amssymb}
\usetheme{default}
\usefonttheme{professionalfonts}
\setbeamercolor{normal text}{fg=black,bg=white}
\setbeamercolor{structure}{fg=black}
\setbeamercolor{frametitle}{fg=black,bg=white}
\setbeamerfont{frametitle}{size=\Large,series=\bfseries}
\setbeamertemplate{navigation symbols}{}
% Default beamer sets the frame title flush to the top-left, outside the text
% block: it started 3mm from the left edge while body text starts at 9mm. Align
% it with the body and give it room to breathe above.
\setbeamertemplate{frametitle}{%
  \vskip5mm\hspace*{2mm}%
  \usebeamerfont{frametitle}\usebeamercolor[fg]{frametitle}\insertframetitle%
  \par\vskip3mm}
\setbeamertemplate{footline}[frame number]
\setbeamertemplate{itemize items}[circle]
\setbeamersize{text margin left=9mm,text margin right=9mm}
\newcommand{\PI}{\mathcal P^+(I)}
\newcommand{\Up}{\mathord{\supseteq}}
\newcommand{\Upstrict}{\mathord{\supsetneq}}
\title{Condensation}
\subtitle{Correspondence}
\author{Satya Benson}
\date{ILIAD, 21 September 2026}
\begin{document}

\begin{frame}{Correspondence between collections}
Two perfect condensations of the same observables satisfy
\[
H(Y_{\Up A}\mid Z_{\Up A})=H(Z_{\Up A}\mid Y_{\Up A})=0.
\]
Here $Y_{\Up A}=(Y_C:A\subseteq C)$ is the latent at $A$ together with those
at all of its supersets. The corresponding collections determine one another
almost surely.
\medskip

The proof has three steps: combine two witnesses, transfer through observations,
and intersect contributing families.
\end{frame}

\begin{frame}{A common joint law}
The two models specify laws for $(X,Y)$ and $(X,Z)$ with the same $X$ marginal.
Choose a joint law for $(X,Y,Z)$ preserving both.
\medskip

One possible choice is
\[
p(x,y,z)=p(x)\,p(y\mid x)\,p(z\mid x).
\]
The theorem holds for any such coupling.
\end{frame}

\begin{frame}{The quantitative theorem}
Fix nonempty $A\subseteq I$ and let $k=|A|$.
For latent models $Y,Z$ of the same observables,
\[
H(Y_{\Up A}\mid Z_{\Up A})
\le \sum_{i\in A}H(Y_{\Up A}\mid X_i)+(k-1)\eta_Z.
\]
Here $\eta_Z$ bounds conditional dependence between upward-closed families
given their intersection. If every reconstruction defect and $\eta_Z$ are at
most $\varepsilon$, the bound reads
\[
H(Y_{\Up A}\mid Z_{\Up A})\le(2k-1)\varepsilon .
\]

This direction uses recovery of the tower from the observations and the
Markov defect of $Z$.
\end{frame}



\begin{frame}{The two-witness lemma}
For finite-valued $U,M,P,Q$,
\[
\begin{aligned}
H(U\mid M)\le{}&H(U\mid M,P)+H(U\mid M,Q)\\
&+I(P;Q\mid M).
\end{aligned}
\]
If both witnesses $(M,P)$ and $(M,Q)$ recover $U$, and their remaining
information is conditionally independent given $M$, then $M$ recovers $U$.
\medskip

The inequality also tracks errors when these requirements hold approximately.
\end{frame}

\begin{frame}{The entropy identity behind the lemma}
Subtract the left side from the right side:
\[
\begin{aligned}
&H(U\mid M,P)+H(U\mid M,Q)+I(P;Q\mid M)\\
&\hspace{25mm}-H(U\mid M)\\[1mm]
&\qquad=I(P;Q\mid U,M)+H(U\mid M,P,Q)\ \ge\ 0.
\end{aligned}
\]
This follows from the chain rule,
\[
H(U,V\mid W)=H(U\mid W)+H(V\mid U,W),
\]
and nonnegativity.
\end{frame}

\begin{frame}{Transfer through an observation}
The contributing family $\Up\{i\}$ satisfies
\[
H(X_i\mid Z_{\Up\{i\}})=0.
\]
Therefore, for any target $U$,
\[
H(U\mid Z_{\Up\{i\}})\le H(U\mid X_i).
\]
The contributing latents determine $X_i$, so they provide at least its
information about $U$.
\end{frame}

\begin{frame}{Intersecting the contributing families}
For $A=\{i_1,\ldots,i_k\}$, set
\[
\mathcal G_j=\bigcap_{\ell=1}^{j}\Up\{i_\ell\}.
\]
An address $C$ survives the final intersection exactly when it contains every
element of $A$:
\[
\mathcal G_k=\{C:A\subseteq C\}=\Up A.
\]
Each $\Up\{i\}$ is upward closed. Intersections preserve upward closure.
\end{frame}

\begin{frame}{The smaller index carries the bigger tuple}
Each intersection adds a requirement, so the families shrink:
\[
\mathcal G_1\supseteq\mathcal G_2\supseteq\cdots\supseteq\mathcal G_k=\Up A.
\]
An address survives step $j$ only if it already survived step $j-1$.
\medskip

A family with more addresses carries a longer tuple of latents, so the tuples
run the other way: every latent in $Z_{\mathcal G_j}$ also appears in
$Z_{\mathcal G_{j-1}}$, and also in $Z_{\Up\{i_j\}}$, since
$\mathcal G_j=\mathcal G_{j-1}\cap\Up\{i_j\}$.
\end{frame}

\begin{frame}{One intersection step}
Use the two-witness lemma with
\[
M=Z_{\mathcal G_j},\qquad
P=Z_{\mathcal G_{j-1}},\qquad
Q=Z_{\Up\{i_j\}}.
\]
The tuple $M$ is already contained in both $P$ and $Q$. Hence
\[
H(U\mid Z_{\mathcal G_j})\le H(U\mid Z_{\mathcal G_{j-1}})
+H(U\mid Z_{\Up\{i_j\}})
+I(Z_{\mathcal G_{j-1}};Z_{\Up\{i_j\}}\mid Z_{\mathcal G_j}).
\]
Both families are upward closed and $\mathcal G_j$ is their intersection, so the
last term is at most $\eta_Z$.
\end{frame}

\begin{frame}{Counting the error terms}
Iterating the previous step down from $\mathcal G_k=\Up A$ to
$\mathcal G_1=\Up\{i_1\}$ uses $k-1$ intersections, and leaves $k$ terms,
one per element of $A$. For any target $U$,
\[
H(U\mid Z_{\Up A})\le\sum_{i\in A}H(U\mid Z_{\Up\{i\}})+(k-1)\eta_Z.
\]
Transfer through an observation bounds each of those $k$ terms by
$H(U\mid X_i)$. Taking the tower $Y_{\Up A}$ as the target turns them into
the reconstruction defects of $Y$:
\[
H(Y_{\Up A}\mid Z_{\Up A})
\le\sum_{i\in A}H(Y_{\Up A}\mid X_i)+(k-1)\eta_Z.
\]
If every reconstruction defect is at most $\delta$ and $\eta_Z\le\eta$, the
right side is $k\delta+(k-1)\eta$.
\end{frame}

\begin{frame}{What the bound compares}
The reverse direction has the same form:
\[
H(Z_{\Up A}\mid Y_{\Up A})
\le\sum_{i\in A}H(Z_{\Up A}\mid X_i)+(k-1)\eta_Y.
\]
For perfect models both sides have zero recovery error and Markov defect,
so both conditional entropies vanish. These are bounds on \textbf{whole collections}.
\medskip

Small errors for each individual latent can accumulate across a collection, so
the tower hypothesis $H(Y_{\Up A}\mid X_i)$ has to be checked directly.
\end{frame}

\begin{frame}{The conditioned score}
For a nonempty observation block $B$, define
\[
\mathcal J_B=\{A:A\cap B\ne\varnothing\},\qquad Y_{\Upstrict A}=(Y_C:A\subsetneq C).
\]
The score charges each active latent given its strict supersets:
\[
\chi(B)=\sum_{A\in\mathcal J_B}H(Y_A\mid Y_{\Upstrict A})
\ge H(Y_{\mathcal J_B})\ge H(X_B).
\]
The chain rule gives the first inequality; LVM gives the second.
\end{frame}

\begin{frame}{Two terms in the score gap}
Let $V_B=\chi(B)-H(Y_{\mathcal J_B})$. Then
\[
\chi(B)-H(X_B)=V_B+H(Y_{\mathcal J_B}\mid X_B).
\]
Both terms are nonnegative.
\begin{itemize}
\item $V_B$ measures departure from the ordered Markov factorization.
\item $H(Y_{\mathcal J_B}\mid X_B)$ measures latent information unrecovered
from the observations.
\end{itemize}
\medskip
\small More precisely, $V_B=D_{\mathrm{KL}}(P_{Y_{\mathcal J_B}}\Vert Q_B)$,
where $Q_B(y)=\prod_{A\in\mathcal J_B}P(y_A\mid y_{\Upstrict A})$.
\end{frame}

\begin{frame}{The score separates models the conditions do not}
Three latent models of the clustering example, all satisfying LVM and the
Markov condition. Take $B$ to be a few points inside one cluster.
\begin{itemize}
\item One latent carrying the whole tuple: $\chi(B)=H(X)$. Every query pays
for all 36 points.
\item A chain, with $X_j$ carried at the address $\{j,\ldots,36\}$:
$\chi(B)=H(X_1,\ldots,X_{\max B})$. Cheap only for early points.
\item The hierarchy, with the scale, the group and the cluster centers above
the points: writing $\Theta$ for the parameters $B$ sits under,
$\chi(B)=H(\Theta)+\sum_{i\in B}H(X_i\mid\Theta)$.
\end{itemize}
Only the hierarchy charges a local query locally, and its gap is
$\chi(B)-H(X_B)=H(\Theta\mid X_B)$, the parameter uncertainty a finite block
leaves. At $B=I$ all three pay $H(X)$: the ranking is relative to the blocks
actually queried.
\end{frame}

\begin{frame}{Why expect agents to satisfy these conditions?}
An agent that answers queries about a block $B$ pays $\chi(B)$, and no
latent model does better than $H(X_B)$. The gap splits into exactly our two
defects:
\[
\chi(B)-H(X_B)=V_B+H(Y_{\mathcal J_B}\mid X_B).
\]
Short descriptions for the queries an agent cares about \emph{are} small
defects. Taking $B=\{i\}$, where $\mathcal J_{\{i\}}=\Up\{i\}$, gives for
$i\in A$
\[
H(Y_{\Up A}\mid X_i)\le H(Y_{\Up\{i\}}\mid X_i)\le\chi(\{i\})-H(X_i),
\]
which is exactly the theorem's reconstruction hypothesis.
\medskip

So the conditions need not be assumed as a coincidence: efficient agents buy
them. Efficiency is relative to the blocks an agent actually queries.
\end{frame}

\end{document}
